\documentclass[11pt,a4paper]{article}
\usepackage[top=2cm,bottom=2.5cm,left=2.5cm,right=2.5cm]{geometry}
\usepackage[french]{babel}
\usepackage{graphicx}
\usepackage[utf8]{inputenc}
\usepackage{verbatim}
\usepackage{fancyhdr}
\usepackage[pdfborder={0 0 0}]{hyperref}
\usepackage{listings}
\lstset{
    basicstyle=\fontsize{8}{10}\selectfont\ttfamily
}
\fancyhead[L]{}
\fancyhead[R]{}
\renewcommand{\footrulewidth}{0.5pt}
\fancyfoot[L]{\textit{Systèmes concurrents}}
\fancyfoot[R]{}
\title{\textbf{Rapport de TP : Systèmes concurrents}}
\author{Charles Follet, Roland Bary}
\date{Fichiers sources disponibles sur \url{https://github.com/cfollet/ada_concurr}}
\begin{document}
\maketitle
\begin{center}
    \maketitle{Master 1 Technologies de l'Internet\\\textbf{Université de Pau et des Pays de l'Adour}}
    \begin{center}
        \includegraphics[scale=0.3]{logoUppa.png}
    \end{center}
\end{center}
\thispagestyle{empty}
\newpage
\pagestyle{fancy}
\renewcommand{\contentsname}{Sommaire}
\tableofcontents
\newpage
\section*{Introduction}
\part{Rendez-vous ADA}
L'objectif ici est de résoudre à l'aide de l'outil RDV ADA les problèmes de:
\begin{itemize}
    \item Producteur/Consommateur
          \begin{itemize}
            \item[•]1 Producteur/ 1Consommateur: tampon de taille N
            \item[•]1 producteur/1 Consommateur avec un tampon de taille 1
          \end{itemize}
    \item Lecteurs/Rédacteurs
          \begin{itemize}
            \item[•] Priorité aux Lecteurs
            \item[•] Priorité aux Rédacteurs
            \item[•] Priorité égales
          \end{itemize}
\end{itemize}
\section{Producteur/ Consommateur}
\subsection{1 Producteur/ 1 Consommateur: tampon de taille N}
\begin{scriptsize}
    \verbatiminput{../prod_cons/rdv_ada/prodcons.adb}
\end{scriptsize}
\newpage
\subsection{Producteur/ Consommateur avec un tampon de taille 1}
\begin{scriptsize}
    \verbatiminput{../prod_cons/rdv_ada/prodcons_un.adb}
\end{scriptsize}
\newpage
\section{Lecteurs/ Rédacteurs}
\subsection{Priorité aux Lecteurs}
\begin{scriptsize}
    \verbatiminput{../lect_red/rdv_ada/prio_lecteur/lecteurredacteur.adb}
\end{scriptsize}
\newpage
\subsection{Priorité aux Rédacteurs}
\begin{scriptsize}
    \verbatiminput{../lect_red/rdv_ada/prio_redacteur/lecteurredacteur.adb}
\end{scriptsize}
\newpage
\subsection{Priorités égales}
\subsubsection{Principe de la solution}
Une solution naïve au problème aurait été de directement utiliser la solution avec les compteurs de Robert vu en cours, c'est à dire : 
\begin{description}
    \item[Lecture] \#act(ECRIRE) = 0
    \item[Ecriture] \#act(ECRIRE) + \#act(LIRE) = 0
\end{description}
Si on utilise ces conditions, le cas suivant ne fonctionne pas :\\~\\
L1 est en cours\\
\indent R1 arrive \\
\indent L2 arrive\\~\\
En effet, rien n’empêche L2 de passer en lecture au détriment de R1.\\
Pour résoudre ce problème il faut que R1 signale son intention d'entrer en SC et nous avons choisis d'utiliser un drapeau.\\
Un lecteur/redacteur désirant aller en SC doit lever le drapeau(drapeau:= 1). Ainsi, un lecteur/redacteur voulant entrer doit disposer de ce drapeau et le problème précédent est résolu.
\subsubsection{Code ADA}
\begin{scriptsize}
    \verbatiminput{../lect_red/rdv_ada/prio_egale/lecteurredacteur.adb}
\end{scriptsize}
\subsubsection{Tests}
\underline{Test 1 :} \\
L1 est en cours\\
\indent R1 arrive \\
\indent L2 arrive\\
\begin{center}
    
    \begin{tabular}{|c|c|c|c|}
        \hline 
        \rule[-1ex]{0pt}{2.5ex} \textbf{Actions} & \textbf{nbLect}       & \textbf{nbRed} & \textbf{drapeau} \\ 
        \hline 
        \rule[-1ex]{0pt}{2.5ex} L1.LIRE          & 1                     & 0              & 0                \\ 
        \hline 
        \rule[-1ex]{0pt}{2.5ex} arrivée de R1   &                       &                &                  \\ 
        \hline 
        \rule[-1ex]{0pt}{2.5ex} R1.barriere      & 1                     & 0              & 1                \\ 
        \hline 
        \rule[-1ex]{0pt}{2.5ex} R1.debutRed      & R1 bloqué a debutRed &                &                  \\ 
        \hline 
        \rule[-1ex]{0pt}{2.5ex} arrivée de L2   &                       &                &                  \\ 
        \hline 
        \rule[-1ex]{0pt}{2.5ex} L2.barriere      & L2 bloqué a barriere &                &                  \\ 
        \hline 
        \rule[-1ex]{0pt}{2.5ex} L1.finLect       & 0                     & 0              & 1                \\ 
        \hline 
        \rule[-1ex]{0pt}{2.5ex} R1.debutRed      & 0                     & 1              & 1                \\ 
        \hline 
        \rule[-1ex]{0pt}{2.5ex} R1.finRed        & 0                     & 0              & 0                \\ 
        \hline 
        \rule[-1ex]{0pt}{2.5ex} L2.debutLect     & 1                     & 0              & 0                \\ 
        \hline 
        \rule[-1ex]{0pt}{2.5ex} L2.finLect       & 0                     & 0              & 0                \\ 
        \hline 
    \end{tabular} 
\end{center}
\newpage

\underline{Test 2 :} \\
R1 est en cours\\
\indent L1 arrive \\
\indent R2 arrive\\
\begin{center}
    \begin{tabular}{|c|c|c|c|}
        \hline 
        \rule[-1ex]{0pt}{2.5ex} \textbf{Actions} & \textbf{nbLect}              & \textbf{nbRed} & \textbf{drapeau} \\ 
        \hline 
        \rule[-1ex]{0pt}{2.5ex} R1.ecrire        & 0                            & 1              & 1                \\ 
        \hline 
        \rule[-1ex]{0pt}{2.5ex} arrivée de L1   &                              &                &                  \\ 
        \hline 
        \rule[-1ex]{0pt}{2.5ex} L1.barriere      & L1 bloqué a barriere (FIFO) &                &                  \\ 
        \hline 
        \rule[-1ex]{0pt}{2.5ex} arrivée de R2   &                              &                &                  \\ 
        \hline 
        \rule[-1ex]{0pt}{2.5ex} R2.barriere      & R2 bloqué a barriere (FIFO) &                &                  \\ 
        \hline 
        \rule[-1ex]{0pt}{2.5ex} R1.finRed        & 0                            & 0              & 0                \\ 
        \hline 
        \rule[-1ex]{0pt}{2.5ex} L1.debutLect     & 1                            & 0              & 0                \\ 
        \hline 
        \rule[-1ex]{0pt}{2.5ex} R2.debutRed      & R2 bloqué a debutRed        &                &                  \\ 
        \hline 
        \rule[-1ex]{0pt}{2.5ex} L1.finLect       & 0                            & 0              & 0                \\ 
        \hline 
        \rule[-1ex]{0pt}{2.5ex} R2.debutRed      & 0                            & 1              & 1                \\ 
        \hline 
        \rule[-1ex]{0pt}{2.5ex} R2.finRed        & 0                            & 0              & 0                \\ 
        \hline 
    \end{tabular} 
\end{center}
\newpage
\part{Objets/Types protégés ADA 95}
L'objectif de cette partie est de :
\begin{itemize}
    \item Rappeler le principe des Objets/Types protégés en ADA 95 : définition, principe, sémantique, différence avec les packages ADA
    \item Comparer les  Objets/types protégés a d'autres outils comme  les sémaphores, RDV ADA, : pouvoir d'expression, difficulté/facilité d'utilisation, etc..
    \item  Implémenter à l'aide Objets/Types protégés:
          \begin{itemize}
            \item[•] Producteur/Consommateur
            \item[•] Lecteurs/ Rédacteurs: sans priorité et  priorité aux lecteurs
            \item[•] Le problème du Carrefour a sens giratoire  
          \end{itemize}
\end{itemize}
%%Beaucoup plus simples que les moniteurs de hoare
\section{Rappel des principes des Objets/Types protégés en ADA 95}
\subsection{Définition}
Un objet protégé définit des données privées ne pouvant être accédées que par les sous-programmes (fonctions, procédures ou entrées) associés à l'objet protégé.
\subsection{Principe}
Les objets protégés permettent l'encapsulation de données privées qui sont accédées par des fonctions (lecture des données), par des procédures (lectures/écritures des données semblables à celles d'un moniteur) et par des entrées (comme une procédure mais avec une mise en attente possible de l'appelant de la procédure)
\textbf{La syntaxe est la suivante}
\begin{verbatim}
protected Nom_deL_Objet_Protege is
	déclarations des fonctions, procédures et entrées de l'objet protégé
private
	déclaration des données privées e l'objet protégé
end Nom_De_L_Objet_Protege;
\end{verbatim}
\subsection{Sémantique}
\subsubsection{Sémantique de base}
\begin{itemize}
\item La partie privée définit les données ou contenu de l'objet protégé.
\item Les fonctions définissent des actions de lecture du contenu de l'objet protégé (interdiction de modifier la valeur de celles-ci)
\item Les procédures et les entrées définissent des actions de lecture et d'écriture des données de l'objet protégé (contrairement à une procédure ou à une fonction, l'appel à une entrée peut être bloquant)
\item Plusieurs lectures peuvent avoir lieu simultanément
\item Une action d'écriture exclut toute autre action (lecture ou écriture)
\end{itemize}
\subsubsection{Sémantique des entrées}
\begin{itemize}
\item[•]Une entrée permet de définir des traitements sous conditions
\begin{itemize}
\item Chaque entrée possède une file d'attente à laquelle est associée une expression booléenne que l'on appelle \textit{garde} ou \textit{barrière}
\item L'entrée n'est "ouverte" que si cette expression (la garde) est vraie
\item Lorsque l'entrée est \textit{fermée} (la garde est fausse), les appels sur cette entrée sont mis en attente (ils ne seront traités que lorsque la garde redeviendra vraie)
\end{itemize}
\item[•]Une tâche, exécutant le code d'une entrée, peut être placée dans la file d'attente d'une autre entrée par l'instruction \textit{requeue}
\item[•]Une tâche en attente sur une entrée interne est plus prioritaire qu'une tâche faisant un nouvel appel à une entrée (ou procédure) de l'objet protégé.
\end{itemize}
\subsection{Différence avec les packages ADA}
Il n'existe pas de différence particulière entre les Objets/Types protégés et les packages ADA. Ces deux notions ont toutes les deux une spécification (avec une partie privée) et un corps.
\newpage
\section{Implémentation à l'aide des Objets/Types protégés}
\subsection{Producteur/ Consommateur}
\begin{scriptsize}
    \verbatiminput{../prod_cons/objet_protege/prodcons.adb}
\end{scriptsize}
\newpage
\subsection{Lecteurs/ Rédacteurs: sans priorité et priorité lecteurs en ADA 95}
\subsubsection{Priorité égale}
\begin{scriptsize}
    \verbatiminput{../lect_red/objet_protege/prio_egale/lecteurredacteur.adb}
\end{scriptsize}
\newpage
\subsubsection{Priorité lecteurs}
\begin{scriptsize}
    \verbatiminput{../lect_red/objet_protege/prio_lecteur/lecteurredacteur.adb}
\end{scriptsize}
\newpage
\subsection{Le problème du carrefour à sens giratoire}
\begin{scriptsize}
    \verbatiminput{../careffour/giratoire.adb}
\end{scriptsize}
\end{document}